\documentclass[UTF8]{ctexart}
\usepackage{amsmath, amssymb}
\usepackage{booktabs}
\usepackage{geometry}
\usepackage{float}
\geometry{a4paper, margin=2.5cm}

\title{数值分析——非线性方程求根上机报告}
\author{物理52~邓昀凯}
\date{2026年10月19日}

\begin{document}
\maketitle

\section{实验目的}

用五种方法求方程
\[
f(x)=x^{3}+2x^{2}+10x-20=0
\]
在 $x_0=1$ 附近的根。准确解为 $x^*=1.368808107\cdots$，要求精度达到
$10^{-7}$。所用方法为：

\begin{enumerate}
  \item 直接迭代 $x_{k+1}=\dfrac{20-2x_k^2-x_k^3}{10}$；
  \item 直接迭代 $x_{k+1}=\sqrt[3]{20-10x_k-2x_k^2}$；
  \item 方法 (1) 的 Steffensen 加速；
  \item 方法 (2) 的 Steffensen 加速；
  \item Newton 法。
\end{enumerate}

\section{实验方法}

\subsection{停止条件与死循环防护}

相邻两次迭代差小于 $\text{tol}=10^{-8}$ 即停止。为防死循环，各方法均设
最大迭代次数，并对非有限值、发散、Steffensen 分母过小等情况提前退出。

\subsection{Steffensen 加速}

对迭代函数 $\varphi(x)$，Steffensen 加速格式为
\[
x_{k+1}=x_k-\frac{(\varphi(x_k)-x_k)^2}
{\varphi(\varphi(x_k))-2\varphi(x_k)+x_k}.
\]

\subsection{Newton 法}

迭代格式为
\[
x_{k+1}=x_k-\frac{f(x_k)}{f'(x_k)},\qquad
f'(x)=3x^2+4x+10.
\]

\section{实验结果}

各方法结果汇总于表~\ref{tab:res}。误差指 $|x-x^*|$。

\begin{table}[H]
\centering
\small
\caption{五种方法的求根结果（$x_0=1$，$\text{tol}=10^{-8}$）}
\label{tab:res}
\begin{tabular}{clccc}
\toprule
序号 & 方法 & 近似根 $x$ & 迭代次数 & 误差 \\
\midrule
(1) & 直接迭代 (1) & $0.548946478$ & 1000 & $8.20\times10^{-1}$ \\
(2) & 直接迭代 (2) & $-3.162277660$ & 1000 & $4.53\times10^{0}$ \\
(3) & Steffensen (1) & $1.368808107821$ & 5 & $8.21\times10^{-10}$ \\
(4) & Steffensen (2) & $1.368808107821$ & 8 & $8.21\times10^{-10}$ \\
(5) & Newton 法 & $1.368808107821$ & 5 & $8.21\times10^{-10}$ \\
\bottomrule
\end{tabular}
\end{table}

其中方法 (1)、(2) 达到最大迭代次数仍未收敛；方法 (3) 因 Steffensen
分母过小提前退出，但此时已满足精度；方法 (4)、(5) 正常收敛。

\section{结果讨论}

\subsection{直接迭代的失效}

方法 (1)、(2) 在根附近均不收敛。原因是它们的迭代函数在 $x^*$ 处的
导数绝对值大于 1：
\[
|\varphi_1'(x^*)|=\left|\frac{-4x^*-3x^{*2}}{10}\right|\approx 1.11>1,
\qquad
|\varphi_2'(x^*)|>1.
\]
不动点迭代的局部收敛条件不满足，迭代点被推离根，最终在别处振荡或发散，
1000 次后误差仍在 $O(1)$ 量级。

\subsection{Steffensen 与 Newton 的高效性}

方法 (3)、(4)、(5) 都在 10 次以内收敛到 $10^{-9}$ 量级误差，远超
$10^{-7}$ 的要求。Steffensen 加速把原本发散的不动点迭代转化为近似二阶
收敛的方法；Newton 法本身是二阶收敛，5 次迭代即达到机器精度附近的误差。

方法 (3) 报出“Steffensen 分母过小”，这是在已经收敛到 $10^{-14}$ 量级
残差后分母接近零所致，属于数值上的正常现象，并非方法失败。

\subsection{精度与效率的比较}

Newton 法迭代次数最少（5 次），且不需要额外选择加速参数，是本问题中
最经济的方法。两种 Steffensen 方法迭代次数分别为 5 和 8，也明显优于
直接迭代。值得注意的是，直接迭代 (1) 虽然发散，但其 Steffensen 加速
版本反而最快，说明加速方法对迭代函数的选取并不十分敏感。

\section{实验总结}

本次实验得到以下结论：

\begin{enumerate}
  \item 不动点迭代是否收敛取决于迭代函数在根处的导数，导数绝对值大于 1
        时直接迭代不可用；
  \item Steffensen 加速能显著改善甚至挽救原本发散的迭代；
  \item Newton 法收敛快、实现简单，是求解单变量非线性方程的首选；
  \item 实际编程中应设置最大迭代次数和发散检测，避免死循环。
\end{enumerate}

注：本次实验部分代码实现与报告排版借助了 AI 工具。

\end{document}